\chapter{Microbiology}

\medterm{Adaptive immune evasion}-- Mechanisms by which pathogens evade the adaptive immune system. Strategies include antigenic variation, which changes surface proteins to avoid antibody recognition. Some viruses downregulate MHC molecules to avoid T cell recognition. Others establish latent infections to avoid immune detection. Understanding adaptive immune evasion is essential for developing new vaccines and immunotherapies. Pathogens that evade adaptive immunity can cause chronic infections.

\medterm{Adenovirus}-- A non-enveloped DNA virus that causes respiratory infections, conjunctivitis, gastroenteritis, and cystitis. Adenoviruses are highly stable in the environment and can cause outbreaks in closed populations. The virus has over 50 serotypes with different tissue tropisms. Diagnosis is made by PCR, antigen detection, or viral culture. There is no specific antiviral treatment, and management is supportive. Severe disseminated infection can occur in immunocompromised patients. Adenoviruses are used as vectors in gene therapy and vaccine development due to their ability to efficiently deliver DNA to target cells.

\medterm{Airborne precautions}-- Infection control measures used for pathogens transmitted through airborne droplet nuclei. Precautions include negative pressure rooms, N95 respirators, and restricted access. Airborne precautions are used for tuberculosis, measles, varicella, and disseminated herpes zoster. The duration of precautions depends on the pathogen and clinical situation. Understanding airborne precautions is essential for preventing transmission of airborne infections in healthcare settings.

\medterm{Antibiotic resistance genes}-- Genes that encode mechanisms conferring resistance to antibiotics. These genes can be located on chromosomes, plasmids, or transposons. Major resistance genes include mecA (methicillin resistance in S. aureus), vanA/vanB (vancomycin resistance in Enterococcus), and blaKPC (carbapenem resistance). Resistance genes can be spread between bacteria through horizontal gene transfer. The acquisition of resistance genes is a major public health concern. Understanding resistance genes is essential for developing effective treatment strategies and infection control policies, and for guiding antimicrobial stewardship programmes.

\medterm{Antibodies}-- Immunoglobulin proteins produced by B cells that bind to specific antigens. There are five classes: IgG (most abundant), IgM (first produced), IgA (mucosal immunity), IgE (allergic responses), and IgD. Antibodies neutralise pathogens, opsonise for phagocytosis, activate complement, and mediate antibody-dependent cellular cytotoxicity. Monoclonal antibodies are engineered antibodies used as therapeutics. Antibody testing is used to diagnose infections and assess immune status. Understanding antibody structure and function is fundamental to immunology, vaccine development, and the design of monoclonal antibody therapies.

\medterm{Antifungal drug classes}-- Major classes of antifungal agents include polyenes (amphotericin B), azoles (fluconazole, voriconazole), echinocandins (caspofungin, micafungin), and allylamines (terbinafine). Polyenes bind to ergosterol in fungal cell membranes. Azoles inhibit ergosterol synthesis. Echinocandins inhibit cell wall synthesis. Allylamines inhibit ergosterol synthesis at a different step. The choice of antifungal depends on the fungus, site of infection, and patient factors. Understanding antifungal drug classes is essential for selecting appropriate therapy, managing resistance, and optimising treatment outcomes in patients with invasive fungal infections.

\medterm{Antifungal prophylaxis}-- The use of antifungal agents to prevent invasive fungal infections in high-risk patients. Common indications include prolonged neutropenia, hematopoietic stem cell transplantation, and advanced HIV infection. Fluconazole, posaconazole, and micafungin are used for prophylaxis. Understanding antifungal prophylaxis is essential for preventing invasive fungal infections in immunocompromised patients. The decision to use prophylaxis should be based on individual risk assessment.

\medterm{Antifungal resistance}-- The ability of fungi to withstand the effects of antifungal agents. Resistance can be intrinsic or acquired. Azole resistance occurs through efflux pumps and target mutations. Echinocandin resistance occurs through FKS gene mutations. Amphotericin B resistance is rare. Candida auris is a multidrug-resistant species causing healthcare-associated outbreaks. Aspergillus fumigatus azole resistance is emerging, partly due to environmental azole use. Understanding antifungal resistance is essential for guiding treatment decisions, implementing antifungal stewardship, and developing new antifungal agents.

\medterm{Antigenic variation}-- The process by which pathogens change their surface proteins to evade immune recognition. Antigenic drift: gradual changes in surface proteins. Antigenic shift: major changes through reassortment of genome segments. Influenza virus demonstrates both mechanisms. Trypanosomes and Plasmodium also use antigenic variation. Understanding antigenic variation is essential for developing effective vaccines and predicting epidemic patterns.

\medterm{Antimicrobial adverse effects}-- Side effects associated with antimicrobial agents. Beta-lactams can cause allergic reactions. Fluoroquinolones can cause tendon rupture and QT prolongation. Aminoglycosides can cause nephrotoxicity and ototoxicity. Azole antifungals can cause hepatotoxicity. Understanding adverse effects is essential for monitoring patients and preventing harm. The benefits of antimicrobial therapy must be weighed against the risks.

\medterm{Antimicrobial dosing optimization}-- Strategies to ensure antimicrobial agents achieve therapeutic concentrations at the site of infection. Optimization includes appropriate drug selection, dosing interval, route of administration, and duration of therapy. Extended or continuous infusion of beta-lactams may improve outcomes. Weight-based dosing is important for many agents. Understanding dosing optimization is essential for maximizing efficacy and minimizing toxicity.

\medterm{Antimicrobial drug interactions}-- Interactions between antimicrobial agents and other medications. Fluoroquinolones chelate divalent cations. Macrolides inhibit CYP3A4. Rifampin is a potent enzyme inducer. Azole antifungals inhibit CYP3A4. Understanding drug interactions is essential for safe prescribing and preventing adverse effects. Careful review of concomitant medications is essential when prescribing antimicrobials.

\medterm{Antimicrobial pharmacodynamics}-- The relationship between antimicrobial drug concentrations and their effect on microorganisms. Time-dependent antibiotics (beta-lactams) are most effective when concentrations remain above the MIC. Concentration-dependent antibiotics (aminoglycosides) are most effective with high peak concentrations. Understanding pharmacodynamics is essential for optimizing dosing regimens and improving clinical outcomes. Pharmacodynamic principles guide antimicrobial dosing strategies.

\medterm{Antimicrobial pharmacokinetics}-- The absorption, distribution, metabolism, and elimination of antimicrobial agents. Pharmacokinetics varies by drug class, patient factors, and disease state. Renal impairment affects elimination of many antibiotics. Hepatic impairment affects metabolism of many antifungals. Understanding pharmacokinetics is essential for dosing optimization in special populations including critically ill patients and those with organ dysfunction.

\medterm{Antimicrobial resistance surveillance}-- The systematic monitoring of antimicrobial resistance patterns in clinical isolates. Surveillance programmes track resistance rates over time and guide empiric therapy. National and international networks share resistance data. Understanding resistance surveillance is essential for effective antimicrobial stewardship and infection control. Surveillance data inform treatment guidelines and public health interventions.

\medterm{Antimicrobial susceptibility testing}-- Laboratory methods used to determine the susceptibility of microorganisms to antimicrobial agents. Methods include disk diffusion, broth microdilution, and automated systems. MIC (minimum inhibitory concentration) is the gold standard measure. Results guide antimicrobial therapy and are interpreted using standardised breakpoints. Resistance testing is essential for managing infections caused by multidrug-resistant organisms. Understanding susceptibility testing is essential for appropriate antimicrobial prescribing and for monitoring resistance trends at local, national, and global levels.

\medterm{Antiretroviral therapy}-- Treatment of HIV infection with combinations of antiretroviral agents. ART targets multiple steps in the viral life cycle including entry, reverse transcription, integration, and maturation. Combination ART suppresses viral replication and restores immune function. ART is lifelong therapy. Understanding ART is essential for managing HIV infection and preventing transmission. ART has transformed HIV from a fatal disease to a chronic manageable condition.

\medterm{Antiseptics}-- Chemical agents applied to living tissue to reduce microbial load. Antiseptics are used for hand hygiene, skin preparation before procedures, and wound care. Common antiseptics include alcohols, chlorhexidine, povidone-iodine, and hydrogen peroxide. Antiseptics have lower efficacy than disinfectants due to toxicity concerns. The choice of antiseptic depends on the application and microbial spectrum. Hand hygiene with alcohol-based rubs is the most effective way to prevent healthcare-associated infections and is the cornerstone of infection control programmes worldwide.

\medterm{Antiviral drug classes}-- Major classes of antiviral agents include nucleoside analogues (acyclovir, tenofovir), protease inhibitors (ritonavir, atazanavir), reverse transcriptase inhibitors (efavirenz, sofosbuvir), and neuraminidase inhibitors (oseltamivir). Each class targets a specific step in the viral replication cycle. Understanding antiviral drug classes is essential for selecting appropriate therapy and managing drug resistance. Combination therapy is often used to prevent resistance.

\medterm{Bacterial adherence}-- The process by which bacteria attach to host cell surfaces, a critical first step in colonisation and infection. Adherence is mediated by adhesins, which are surface structures that bind to specific receptors on host cells. Adhesins include pili, fimbriae, outer membrane proteins, and capsules. The specificity of adhesin-receptor interactions determines the tissue tropism of bacterial pathogens. For example, uropathogenic E. coli express adhesins that bind to uroepithelial cells. Anti-adhesion strategies, including competitive inhibitors of adhesin-receptor binding, are being investigated as novel therapeutic approaches to prevent bacterial colonisation.

\medterm{Bacterial biofilms}-- Structured communities of bacteria attached to surfaces and embedded in a self-produced matrix of extracellular polymeric substances. Biofilms are a major cause of chronic infections, particularly on medical devices including catheters, prosthetic joints, and heart valves. Bacteria in biofilms are more resistant to antibiotics (up to 1000-fold) and host immune responses. The biofilm matrix consists of polysaccharides, proteins, and extracellular DNA. Quorum sensing regulates biofilm formation and dispersal, making it a potential target for anti-biofilm therapies.

\medterm{Bacterial capsule}-- A polysaccharide layer surrounding the bacterial cell wall that protects against phagocytosis and complement-mediated killing. Capsules are major virulence factors for many pathogens including Streptococcus pneumoniae, Haemophilus influenzae, and Klebsiella pneumoniae. The capsule inhibits opsonisation and prevents engulfment by immune cells. Anti-capsular antibodies are protective and form the basis of conjugate vaccines. Capsule production can be demonstrated by the quellung reaction, which shows capsular swelling when specific antibodies bind to the capsule.

\medterm{Bacterial cell wall}-- The rigid outer layer of bacteria that provides structural integrity and determines cell shape. The cell wall is composed primarily of peptidoglycan, a polymer of sugars cross-linked by short peptides. Gram-positive bacteria have a thick peptidoglycan layer (20-80 nm) while gram-negative bacteria have a thin layer (5-10 nm) surrounded by an outer membrane. The cell wall is essential for bacterial survival and is the target of many antibiotics including beta-lactams and vancomycin. The unique structure of the bacterial cell wall, particularly peptidoglycan synthesis enzymes, provides an excellent target for antimicrobial drugs that have minimal toxicity to human cells.

\medterm{Bacterial classification}-- The organisation of bacteria into taxonomic groups based on genetic, phenotypic, and evolutionary relationships. Bacteria are classified using a hierarchical system: domain, phylum, class, order, family, genus, and species. Classification methods include phenotypic characterisation (morphology, staining, biochemical tests) and genotypic methods (16S rRNA sequencing, DNA-DNA hybridisation). The 16S rRNA gene is the primary marker for bacterial phylogenetics. Modern classification emphasises molecular phylogenetic approaches, particularly whole-genome sequencing, which provides more accurate taxonomic resolution than traditional methods.

\medterm{Bacterial conjugation}-- The transfer of genetic material between bacteria through direct cell-to-cell contact via a pilus. Conjugation is mediated by conjugative plasmids that encode the genes for transfer. The donor cell (F+) forms a conjugative pilus that attaches to the recipient cell (F-). A single strand of plasmid DNA is transferred through the pilus, and both cells synthesise the complementary strand. Conjugation is the most efficient mechanism of horizontal gene transfer and can spread antibiotic resistance genes across bacterial species and genera, contributing to the rapid emergence of multidrug-resistant pathogens.

\medterm{Bacterial dormancy}-- A state of reduced metabolic activity that allows bacteria to survive unfavourable conditions. Dormant forms include spores, persisters, and viable but non-culturable cells. Spores are highly resistant structures formed by Bacillus and Clostridium. Persisters are phenotypic variants that are tolerant to antibiotics. Understanding dormancy is essential for treating chronic infections and developing new antimicrobial strategies.

\medterm{Bacterial endospores}-- Highly resistant dormant structures formed by certain gram-positive bacteria including Bacillus and Clostridium species. Endospores form during nutrient deprivation and can survive for decades. They are resistant to heat, radiation, chemicals, and desiccation. Sporulation involves asymmetric cell division, engulfment of the forespore, cortex formation, and coat assembly. Germination occurs when conditions improve, with the spore returning to vegetative growth. Endospore formation is a significant challenge for infection control because spores resist standard disinfection methods and can survive for extended periods in the healthcare environment.

\medterm{Bacterial flagella}-- Thread-like appendages composed of flagellin protein that provide motility to bacteria. Flagella are composed of three parts: the filament, hook, and basal body. The basal body is embedded in the cell wall and functions as a rotary motor. Bacteria are classified by flagellar arrangement: monotrichous (single flagellum), lophotrichous (tuft at one pole), amphitrichous (flagella at both poles), and peritrichous (flagella distributed over entire surface). Flagella are important for chemotaxis, allowing bacteria to move towards nutrients and away from harmful substances, and they also contribute to pathogenesis by enabling bacteria to reach host tissues.

\medterm{Bacterial genetics}-- The study of genes and genetic processes in bacteria. Bacteria have a single circular chromosome and may contain plasmids, which are extrachromosomal DNA elements. Bacterial genes are organized into operons, which are clusters of genes regulated as a unit. Genetic variation occurs through mutation, transformation, transduction, and conjugation. Horizontal gene transfer allows rapid spread of antibiotic resistance and virulence factors. The bacterial genome is compact with little non-coding DNA, and its analysis provides insights into metabolic capabilities, virulence potential, and evolutionary relationships between strains.

\medterm{Bacterial growth phases}-- The stages of bacterial growth in culture. The lag phase is when bacteria adapt to their environment. The log (exponential) phase is when bacteria multiply at their maximum rate. The stationary phase is when nutrient depletion and waste accumulation limit growth. The death phase is when bacteria die. Understanding growth phases is essential for antimicrobial susceptibility testing and understanding pathogenesis. Many antibiotics are most effective during the log phase.

\medterm{Bacterial metabolism}-- The chemical processes that occur within bacteria to maintain life. Bacteria have diverse metabolic capabilities including aerobic and anaerobic respiration, fermentation, and photosynthesis. Metabolic pathways are targets for antimicrobial drugs. Understanding bacterial metabolism is essential for developing new antibiotics and understanding pathogenesis. Metabolic flexibility allows bacteria to survive in diverse environments.

\medterm{Bacterial morphology}-- The shape and arrangement of bacteria as observed under microscopy. Basic shapes include cocci (spheres), bacilli (rods), and spirilla (spirals). Cocci can be arranged in pairs (diplococci), chains (streptococci), clusters (staphylococci), or tetrads. Bacilli can be single, paired, or arranged in chains. Spiral bacteria include vibrios (curved), spirilla (rigid spirals), and spirochetes (flexible spirals). Morphology is a fundamental characteristic for bacterial identification. The size of bacteria typically ranges from 0.5 to 5 micrometres, making them significantly smaller than eukaryotic cells.

\medterm{Bacterial pathogenesis}-- The mechanisms by which bacteria cause disease. Pathogenesis involves adhesion, invasion, toxin production, and immune evasion. Adhesins mediate attachment to host cells. Invasins allow penetration of tissues. Toxins damage host cells. Immune evasion mechanisms include capsules, biofilms, and intracellular survival. Understanding pathogenesis is essential for developing new antimicrobial therapies and vaccines. The study of pathogenesis has identified virulence factors that are targets for drug development, including anti-virulence therapies that disarm pathogens without killing them, potentially reducing selective pressure for resistance.

\medterm{Bacterial pili}-- Hair-like appendages on the bacterial surface composed of pilin protein. Pili are shorter and thicker than flagella. Type I pili mediate adherence to host cells and are important virulence factors. Type IV pili are involved in twitching motility and DNA uptake during transformation. F-pili (sex pili) facilitate conjugation between bacteria. Pili are antigenic (K antigens) and can stimulate immune responses. They are too thin to be visible with light microscopy. Pili are essential for colonisation and biofilm formation, making them important targets for vaccine development and anti-adhesion therapies.

\medterm{Bacterial toxins}-- Substances produced by bacteria that cause damage to host cells. There are two major categories: exotoxins and endotoxins. Exotoxins are proteins secreted by bacteria that have specific toxic effects. Endotoxins are lipopolysaccharide components of gram-negative bacterial outer membranes released during cell lysis. Exotoxins are classified into types based on mechanism of action: A-B toxins, membrane-disrupting toxins, and superantigens. Endotoxin causes fever, hypotension, and disseminated intravascular coagulation through activation of the innate immune system, and is a key mediator of septic shock in gram-negative bacterial infections.

\medterm{Bacterial transduction}-- The transfer of bacterial DNA from one bacterium to another through bacteriophages (bacterial viruses). There are two types: generalized transduction, which occurs during the lytic cycle when random bacterial DNA is packaged into phage particles, and specialized transduction, which occurs during the lysogenic cycle when specific genes adjacent to the prophage are transferred. Transduction is an important mechanism of horizontal gene transfer and can spread antibiotic resistance genes and virulence factors between bacterial strains, contributing to the evolution of pathogenic bacteria.

\medterm{Bacterial transformation}-- The uptake of free DNA from the environment by competent bacteria. Transformation occurs naturally in some bacterial species including Streptococcus pneumoniae, Haemophilus influenzae, and Neisseria. Competence is the ability to take up DNA and is regulated by specific genes. The DNA can be incorporated into the bacterial chromosome through homologous recombination. Transformation is an important mechanism of horizontal gene transfer and can spread antibiotic resistance genes. Artificial transformation is widely used in molecular biology to introduce engineered plasmids into bacteria for research and industrial applications.

\medterm{Bacterial virulence factors}-- Molecules produced by bacteria that contribute to their ability to cause disease. Virulence factors include adhesins, invasins, toxins, capsules, and enzymes. Adhesins mediate attachment to host cells. Toxins damage host cells directly. Capsules protect against phagocytosis. Enzymes degrade host tissues. Understanding virulence factors is essential for developing new antimicrobial therapies and vaccines. Virulence factors are potential targets for drug and vaccine development.

\medterm{Bactericide}-- An agent that kills bacteria rather than merely inhibiting their growth (bacteriostatic). Bactericidal antibiotics include beta-lactams (penicillins, cephalosporins, carbapenems), aminoglycosides, fluoroquinolones, metronidazole, and vancomycin. The distinction is concentration-dependent and organism-specific; some drugs are bactericidal for certain organisms and bacteriostatic for others. Bactericidal activity is preferred for endocarditis, meningitis, and neutropenic fever.

\medterm{Bacteriology}-- The branch of microbiology concerned with the study of bacteria, including their morphology, genetics, physiology, ecology, and clinical significance. Encompasses bacterial taxonomy, identification methods (Gram stain, culture, biochemical tests, MALDI-TOF, 16S rRNA sequencing), pathogenesis mechanisms, antibiotic susceptibility testing, and the role of bacteria in health (microbiome) and disease. Foundational to clinical infectious disease practice and infection control.

\medterm{Bacteriophages}-- Viruses that specifically infect and replicate within bacteria, also known as phages. Each phage type typically infects a narrow range of bacterial strains. Phages attach to bacterial surface receptors and inject their genome, hijacking bacterial machinery to produce progeny phages that lyse the host cell. Phage therapy, using lytic phages to treat multidrug-resistant bacterial infections, is a re-emerging antimicrobial strategy. Phages are also essential tools in molecular biology (lambda, M13, T4).

\medterm{Biofilm-associated infections}-- Infections caused by bacteria growing in biofilms on medical devices or tissues. Biofilms are structured communities of bacteria embedded in extracellular polymeric substances. Bacteria in biofilms are more resistant to antibiotics and host immune responses. Common biofilm-associated infections include catheter-associated urinary tract infections, ventilator-associated pneumonia, and prosthetic joint infections. Prevention through anti-biofilm strategies and device management is essential. Understanding biofilm-associated infections is essential for developing effective prevention and treatment strategies, including antimicrobial lock therapy and biofilm-disrupting agents.

\medterm{C. difficile}-- Clostridioides difficile, a gram-positive, spore-forming bacterium that causes antibiotic-associated diarrhea and colitis. C. difficile produces toxins A and B that cause colonic inflammation. Risk factors include antibiotic exposure, proton pump inhibitor use, and advanced age. Treatment includes vancomycin or fidaxomicin. Recurrence occurs in approximately 25\% of cases. fecal microbiota transplantation is highly effective for multiply recurrent infections. Prevention through antibiotic stewardship, environmental cleaning with sporicidal agents, and isolation of affected patients is essential to control C. difficile transmission in healthcare settings.

\medterm{Chikungunya virus}-- A positive-sense RNA virus belonging to the Togaviridae family that causes fever and severe joint pain. The virus is transmitted by Aedes mosquitoes. Symptoms include fever, rash, and polyarthralgia that can be debilitating. Most patients recover fully, but some develop chronic joint pain lasting months to years. Diagnosis is made by PCR or serology. Treatment is supportive with pain management. There is no specific antiviral treatment. Prevention through mosquito control is important. The virus continues to cause outbreaks in tropical and subtropical regions where vector control is inadequate, and travellers returning from endemic areas may import the disease.

\medterm{Clostridioides difficile pathogenesis}-- The mechanisms by which C. difficile causes disease. The organism produces toxins A and B that damage the colonic mucosa. Antibiotic disruption of normal gut flora allows C. difficile to colonize and produce toxins. The organism forms spores that are resistant to disinfectants and can persist in the environment. Understanding C. difficile pathogenesis is essential for developing new treatments and prevention strategies. The organism is a major cause of healthcare-associated diarrhea.

\medterm{Clostridium botulinum}-- A gram-positive, spore-forming, anaerobic rod that produces botulinum toxin, the most potent toxin known. The toxin causes flaccid paralysis by blocking acetylcholine release at neuromuscular junctions. Infant botulism occurs from ingestion of spores in honey or soil. Wound botulism occurs from injection of contaminated heroin. Food-borne botulism occurs from ingestion of preformed toxin in improperly preserved food. Diagnosis is clinical with toxin detection. Treatment includes botulinum antitoxin to neutralise circulating toxin and intensive supportive care, including mechanical ventilation if respiratory muscle paralysis occurs.

\medterm{Clostridium perfringens}-- A gram-positive, spore-forming, anaerobic rod that causes gas gangrene (clostridial myonecrosis) and food poisoning. The organism produces alpha-toxin, which causes tissue necrosis. Gas gangrene presents with crepitus, tissue necrosis, and systemic toxicity. Food poisoning presents with acute onset of diarrhea and abdominal cramps. Diagnosis involves culture and toxin detection. Treatment includes surgical debridement, hyperbaric oxygen, and penicillin. The organism is part of the normal gastrointestinal flora in some individuals, but when introduced into deep tissues through trauma or surgery, it can cause rapidly progressive necrotising infection.

\medterm{Clostridium tetani}-- A gram-positive, spore-forming, anaerobic rod that causes tetanus. The organism produces tetanospasmin, a neurotoxin that causes spastic paralysis. Spores are found in soil and enter through wound contamination. The toxin blocks inhibitory neurotransmitters causing uncontrolled muscle contractions. Symptoms include trismus (lockjaw), opisthotonus, and respiratory failure. Diagnosis is clinical. Treatment includes metronidazole, tetanus immune globulin, and supportive care. Prevention through vaccination with tetanus toxoid and appropriate wound management is highly effective.

\medterm{Contact precautions}-- Infection control measures used for patients colonised or infected with multidrug-resistant organisms. Precautions include gown and glove use, patient placement in single rooms, and dedicated equipment. Contact precautions are used for MRSA, VRE, CRE, and other resistant organisms. Hand hygiene remains essential in addition to contact precautions. The duration of precautions depends on the organism and clinical situation. Understanding contact precautions is essential for preventing transmission of multidrug-resistant organisms in healthcare settings and reducing the burden of healthcare-associated infections.

\medterm{Coxiella burnetii}-- An obligate intracellular bacterium that causes Q fever, a zoonotic infection transmitted from animals to humans. The organism is transmitted through inhalation of contaminated aerosols or ingestion of contaminated dairy products. Acute Q fever presents with fever, headache, and pneumonia. Chronic Q fever can cause endocarditis. Diagnosis involves serology and PCR. Treatment includes doxycycline for acute disease and doxycycline with hydroxychloroquine for chronic endocarditis. Prevention through pasteurisation of dairy products, protective measures for abattoir workers and veterinarians, and awareness of the disease in endemic regions is essential for reducing transmission.

\medterm{CRE}-- Carbapenem-resistant Enterobacteriaceae, resistant to carbapenem antibiotics through production of carbapenemases including KPC, NDM, and OXA-48-like enzymes. CRE are a major public health threat due to limited treatment options. Infections are associated with high mortality rates. Treatment options include ceftazidime-avibactam, meropenem-vaborbactam, and aminoglycosides. Infection prevention and antimicrobial stewardship are essential to control spread. Active surveillance may be necessary in high-risk settings such as intensive care units to detect colonised patients and implement timely infection control measures.

\medterm{Cytomegalovirus}-- A herpesvirus that causes significant disease in immunocompromised patients including transplant recipients and HIV-infected individuals. CMV infects a wide range of cell types and establishes latent infections in myeloid progenitor cells. Primary infection is usually asymptomatic in immunocompetent individuals. Reactivation can cause retinitis, colitis, oesophagitis, and pneumonitis in immunocompromised patients. Congenital CMV infection is the most common viral cause of birth defects. Diagnosis involves PCR, antigen detection, viral culture, and serology, with treatment including ganciclovir, valganciclovir, or foscarnet for active disease in immunocompromised patients.

\medterm{Dengue virus}-- A positive-sense RNA virus belonging to the Flaviviridae family that causes dengue fever. The virus has four serotypes, and infection with one serotype provides lifelong immunity to that serotype but only temporary cross-protection against others. Secondary infection with a different serotype can cause severe dengue hemorrhagic fever or dengue shock syndrome due to antibody-dependent enhancement. The virus is transmitted by Aedes mosquitoes. Diagnosis is made by PCR or antigen detection. Treatment is supportive with careful fluid management and avoidance of non-steroidal anti-inflammatory drugs due to the risk of bleeding complications.

\medterm{Diagnostic microbiology}-- The laboratory identification of microorganisms causing disease. Methods include microscopy, culture, biochemical testing, serology, and molecular methods. Rapid diagnostic tests provide same-day results. MALDI-TOF mass spectrometry revolutionized bacterial identification. PCR and nucleic acid amplification tests are essential for diagnosing infections caused by fastidious organisms. Understanding diagnostic microbiology is essential for appropriate antimicrobial therapy.

\medterm{Disinfection}-- The elimination of most pathogenic microorganisms on inanimate surfaces. Disinfection does not eliminate all forms of microbial life, particularly spores. High-level disinfection eliminates all microorganisms except spores. Intermediate-level disinfection eliminates mycobacteria, most viruses, and bacteria. Low-level disinfection eliminates most bacteria and some viruses. Chemical disinfectants include alcohols, bleach, quaternary ammonium compounds, and phenolics. The level of disinfection required depends on the intended use of the medical device, with critical items requiring sterilisation, semi-critical items requiring high-level disinfection, and non-critical items requiring intermediate or low-level disinfection.

\medterm{Droplet precautions}-- Infection control measures used for pathogens transmitted through respiratory droplets. Precautions include surgical masks, patient placement in single rooms, and spatial separation. Droplet precautions are used for influenza, pertussis, and meningococcal infections. The distance for spatial separation is typically 3-6 feet. Understanding droplet precautions is essential for preventing transmission of respiratory infections in healthcare settings.

\medterm{Ebola virus}-- A filovirus that causes Ebola virus disease, a severe hemorrhagic fever. The virus is transmitted through contact with infected body fluids. The incubation period is 2-21 days. Symptoms include fever, headache, myalgia, and hemorrhage. Diagnosis is made by PCR and antigen detection. Treatment is supportive with rehydration and management of complications. The rVSV-ZEBOV vaccine is available. Outbreaks have occurred primarily in Africa with significant public health impact. Prevention through infection control measures, including isolation of affected patients, use of personal protective equipment, and safe burial practices, is critical to containing Ebola outbreaks.

\medterm{Echinococcus granulosus}-- A tapeworm that causes hydatid disease (cystic echinococcosis). The adult worm infects dogs, and humans are accidental intermediate hosts. Infection occurs through ingestion of eggs from dog feces. Larval cysts form in the liver and lungs, causing space-occupying lesions. Diagnosis involves imaging showing characteristic cysts and serology. Treatment includes surgical removal or percutaneous aspiration with albendazole coverage. Puncture of cysts can cause anaphylaxis. Prevention through regular deworming of dogs and proper disposal of dog feces is essential in endemic areas to break the transmission cycle.

\medterm{Ehrlichia and Anaplasma}-- Obligate intracellular bacteria transmitted by ticks that cause ehrlichiosis and anaplasmosis. Ehrlichia chaffeensis causes human monocytic ehrlichiosis. Anaplasma phagocytophilum causes human granulocytic anaplasmosis. The organisms infect white blood cells (Ehrlichia) or granulocytes (Anaplasma). Symptoms include fever, headache, and myalgia. Diagnosis involves PCR and serology. Treatment includes doxycycline. Prevention through tick avoidance is essential. These infections can be severe in immunocompromised patients and may require prolonged antibiotic therapy to prevent relapse.

\medterm{Endemic fungi}-- Fungi that are found in specific geographic regions and cause systemic infections in immunocompetent individuals. Endemic fungi include Histoplasma capsulatum (Ohio and Mississippi River valleys), Blastomyces dermatitidis (Great Lakes region), Coccidioides immitis (southwestern United States), and Paracoccidioides brasiliensis (Central and South America). These fungi are dimorphic, existing as moulds in the environment and yeasts in tissue. Infections are acquired through inhalation of spores. Diagnosis involves culture, histopathology, and antigen detection in tissue specimens, with treatment including azole antifungals or amphotericin B depending on the species and severity.

\medterm{Endotoxins}-- Lipopolysaccharide components of gram-negative bacterial outer membranes released during cell lysis or bacterial growth. Endotoxin is composed of lipid A (the toxic component), core oligosaccharide, and O-antigen. It activates macrophages through toll-like receptor 4, leading to cytokine release. Clinical effects include fever, hypotension, and disseminated intravascular coagulation. Endotoxin is heat-stable and difficult to remove. All gram-negative bacteria contain endotoxin, but potency varies between species, with Neisseria meningitidis and Escherichia coli producing particularly potent endotoxin.

\medterm{Entamoeba histolytica}-- A protozoan parasite that causes amoebiasis, a gastrointestinal and extra-intestinal infection transmitted via the fecal-oral route. The organism exists in two forms: the infectious cyst (transmitted in contaminated food or water) and the invasive trophozoite. Trophozoites colonise the large intestine, where they can adhere to colonic epithelium via a galactose-binding lectin and cause tissue invasion, often producing characteristic flask-shaped ulcers. In most cases, infection is asymptomatic; however, invasive disease can cause amoebic dysentery (bloody diarrhea with mucus, abdominal pain, and tenesmus). Extraintestinal spread most commonly results in amoebic liver abscess. Diagnosis involves microscopy demonstrating trophozoites with ingested red blood cells, antigen detection, PCR, and serology for extra-intestinal disease. Treatment includes tinidazole or metronidazole followed by a luminal agent (paromomycin or diloxanide furoate) to eradicate cyst carriage.

\medterm{Epstein-Barr virus}-- A herpesvirus that causes infectious mononucleosis and is associated with several malignancies including Burkitt lymphoma, nasopharyngeal carcinoma, and Hodgkin lymphoma. EBV infects B cells and establishes latent infections. The virus encodes several oncoproteins including LMP1 and EBNA2. Primary infection is often asymptomatic in childhood but can cause mononucleosis in adolescents and adults. Diagnosis is confirmed by heterophile antibody test or EBV-specific serology. There is no specific antiviral treatment for EBV, and management of infectious mononucleosis is supportive, with corticosteroids reserved for airway compromise or severe complications.

\medterm{ESBL}-- Extended-spectrum beta-lactamase, enzymes that hydrolyse penicillins, cephalosporins, and monobactams. ESBLs are common in Enterobacteriaceae, particularly Klebsiella pneumoniae and E. coli. ESBL production confers resistance to most beta-lactam antibiotics except carbapenems. Infections are associated with increased morbidity and mortality. Treatment typically requires carbapenems or other agents active against ESBL-producing organisms. Stewardship strategies are essential to limit the spread of ESBL-producing organisms and preserve the efficacy of remaining treatment options.

\medterm{Escherichia coli}-- A gram-negative rod that is part of the normal intestinal flora but can cause various infections. Uropathogenic E. coli causes urinary tract infections. Enterotoxigenic E. coli causes traveller's diarrhea. Enterohemorrhagic E. coli (O157:H7) causes bloody diarrhea and hemolytic uraemic syndrome. ESBL-producing E. coli is an increasing cause of healthcare-associated infections. The organism is identified by lactose fermentation and biochemical tests. Treatment depends on the infection site and susceptibility profile, with increasing resistance to commonly used antibiotics necessitating tailored therapy based on local resistance patterns.

\medterm{Exotoxins}-- Proteins secreted by bacteria that have specific toxic effects on host cells. Exotoxins are generally more potent than endotoxins and have specific mechanisms of action. A-B toxins consist of an active (A) component and a binding (B) component. Examples include diphtheria toxin, cholera toxin, and botulinum toxin. Membrane-disrupting toxins include haemolysins and leukocidins that damage cell membranes. Superantigens including toxic shock syndrome toxin-1 cause non-specific T cell activation. Exotoxins are a major focus of vaccine development, as demonstrated by toxoid vaccines for diphtheria and tetanus, and they are also used in biomedical research for their specific cellular targets.

\medterm{Extracellular pathogens}-- Microorganisms that survive and replicate outside host cells. Extracellular bacteria include Streptococcus, Staphylococcus, and Escherichia coli. These pathogens are typically killed by antibodies and complement. They produce toxins and enzymes to damage host tissues. Understanding extracellular pathogenesis is essential for developing vaccines that generate protective antibodies. Extracellular pathogens are typically susceptible to antibiotics that can reach high concentrations in extracellular fluids and are effectively cleared by antibody-mediated immune mechanisms.

\medterm{Filariasis}-- A group of nematode infections transmitted by mosquitoes and other vectors. The major species causing lymphatic filariasis are Wuchereria bancrofti, Brugia malayi, and Brugia timori. Adult worms reside in lymphatic vessels causing lymphoedema and elephantiasis. Diagnosis involves microscopy of blood collected at night when microfilariae are circulating. Treatment includes diethylcarbamazine (DEC). Prevention through mosquito control is essential. Onchocerciasis (river blindness) is caused by Onchocerca volvulus, transmitted by black flies, and is treated with ivermectin.

\medterm{Filoviruses}-- A family of negative-sense RNA viruses that includes Ebola and Marburg viruses. These viruses cause severe hemorrhagic fever with high mortality rates. The virus is transmitted through contact with infected body fluids. The incubation period is 2-21 days. Symptoms include fever, headache, myalgia, and hemorrhage. Diagnosis is made by PCR and antigen detection. Treatment is supportive with rehydration and management of complications. The rVSV-ZEBOV vaccine is available for Ebola virus. Prevention through infection control measures, including isolation and barrier nursing, is essential for managing filovirus outbreaks in healthcare settings.

\medterm{Francisella tularensis}-- A gram-negative coccobacillus that causes tularemia, a zoonotic infection transmitted by ticks, deer flies, or direct contact with infected animals. The organism can be acquired through the skin, respiratory tract, or gastrointestinal tract. Ulceroglandular tularemia presents with skin ulcer and regional lymphadenopathy. Pneumonic tularemia can be severe. Diagnosis involves serology and culture (requires special media). Treatment includes streptomycin or gentamicin. Tularemia is a rare infection that is considered a potential bioterrorism agent due to its low infectious dose and ease of aerosolisation.

\medterm{Fungal pathogenesis}-- The mechanisms by which fungi cause disease. Fungal pathogenesis involves adhesion, invasion, and evasion of host defences. Adhesins mediate attachment to host cells. Invasive hyphae penetrate tissues. Capsules and melanin protect against immune responses. Biofilm formation provides protection against antifungal agents. Immunocompromised individuals are at increased risk for invasive fungal infections. Understanding fungal pathogenesis is essential for developing new antifungal therapies and vaccines, particularly given the increasing population of immunocompromised patients at risk for invasive fungal disease.

\medterm{Germ}-- A colloquial term for a microorganism that causes disease.

\medterm{Giardia lamblia}-- A protozoan parasite that causes giardiasis, a diarrhoeal illness transmitted through the fecal-oral route. The organism has a trophozoite and cyst stage. Infection occurs through ingestion of cysts in contaminated water. Symptoms include foul-smelling diarrhea, abdominal cramps, and bloating. Diagnosis involves stool microscopy showing trophozoites or cysts, and antigen detection. Treatment includes metronidazole or tinidazole. Prevention through water treatment and sanitation is essential. Giardia lamblia is a common cause of waterborne outbreaks and is highly prevalent in areas with poor sanitation and inadequate water treatment.

\medterm{Gram stain}-- A differential staining technique that classifies bacteria into two major groups based on cell wall structure. The procedure involves crystal violet staining, iodine mordant, alcohol decolorisation, and safranin counterstain. Gram-positive bacteria retain the crystal violet-iodine complex due to their thick peptidoglycan layer and appear purple. Gram-negative bacteria lose the crystal violet and take up the counterstain, appearing pink or red. The Gram stain is the most important initial step in bacterial identification and provides critical information for guiding empiric antimicrobial therapy while cultures are pending.

\medterm{Haemophilus influenzae}-- A gram-negative coccobacillus that causes meningitis, epiglottitis, pneumonia, and other infections. The encapsulated type b (Hib) strain was the most virulent, causing invasive disease in children. The Hib conjugate vaccine has dramatically reduced the incidence of invasive Hib disease. Non-typeable strains cause otitis media, sinusitis, and exacerbations of chronic bronchitis. The organism requires X factor (haemin) and V factor (NAD) for growth. Diagnosis involves culture on chocolate agar. Treatment includes broad-spectrum antibiotics such as third-generation cephalosporins, and prevention of invasive Hib disease has been achieved through routine childhood vaccination.

\medterm{Hand hygiene}-- The most effective measure for preventing healthcare-associated infections. Hand hygiene includes handwashing with soap and water and hand rubbing with alcohol-based rubs. Alcohol-based rubs are preferred for routine hand hygiene as they are faster and more effective than soap and water. Soap and water should be used when hands are visibly soiled or when caring for patients with C. difficile. The WHO has established five moments for hand hygiene in healthcare settings. Hand hygiene compliance is monitored in healthcare settings as a key quality indicator, and improvement strategies include education, feedback, and the availability of alcohol-based hand rub at the point of care.

\medterm{Helminths}-- Parasitic worms that cause disease in humans. The major groups are nematodes (roundworms), cestodes (tapeworms), and trematodes (flukes). Helminths have complex life cycles often involving intermediate hosts. Infection occurs through ingestion of eggs or larvae, penetration of skin, or insect vectors. Diagnosis involves stool microscopy for eggs and larvae, serology, and imaging. Treatment includes anthelmintic agents such as albendazole, mebendazole, and praziquantel. Helminth infections are a major cause of morbidity in tropical and subtropical regions, contributing to malnutrition, anemia, and impaired cognitive development in children.

\medterm{Herpes simplex virus}-- A double-stranded DNA virus belonging to the Herpesviridae family that causes oral and genital herpes. HSV has two serotypes: HSV-1, which typically causes oral herpes, and HSV-2, which typically causes genital herpes. The virus establishes latent infections in sensory ganglia and can reactivate periodically. Primary infection is often asymptomatic. Reactivation causes recurrent outbreaks of vesicular lesions. Antiviral treatment with acyclovir, valacyclovir, or famciclovir can reduce the frequency and severity of recurrent episodes and is recommended for patients with frequent or severe outbreaks.

\medterm{Hookworms}-- Nematodes including Ancylostoma duodenale and Necator americanus that cause hookworm infection. Infection occurs through skin penetration by filariform larvae, typically when walking barefoot on contaminated soil. Adult worms attach to the intestinal mucosa and feed on blood, causing iron deficiency anemia. Diagnosis involves stool microscopy for eggs. Treatment includes albendazole or mebendazole. Prevention through sanitation, wearing shoes, and avoiding soil contact is essential. Hookworm infection remains a significant public health problem in resource-limited settings, particularly among children and agricultural workers in tropical regions.

\medterm{Horizontal gene transfer}-- The transfer of genetic material between bacteria through mechanisms other than parent-to-offspring transmission. The three mechanisms are transformation (uptake of free DNA), transduction (transfer via bacteriophages), and conjugation (direct cell-to-cell transfer). Horizontal gene transfer allows rapid spread of antibiotic resistance genes and virulence factors. Understanding this process is essential for tracking the spread of resistance and developing strategies to combat it.

\medterm{Hospital-acquired infections}-- Infections that develop during or after hospitalisation. Major types include central line-associated bloodstream infections, catheter-associated urinary tract infections, ventilator-associated pneumonia, and surgical site infections. Risk factors include invasive devices, prolonged hospitalisation, and antimicrobial use. Prevention through infection control measures including hand hygiene, device bundles, and antimicrobial stewardship is essential. Hospital-acquired infections are associated with increased morbidity, mortality, length of stay, and healthcare costs, and represent a major patient safety concern.

\medterm{Human immunodeficiency virus}-- A retrovirus belonging to the Lentivirus genus that causes acquired immunodeficiency syndrome (AIDS). HIV has two species: HIV-1 and HIV-2. The virus targets CD4+ T cells, macrophages, and dendritic cells. HIV has a complex genome encoding structural proteins, enzymes (reverse transcriptase, integrase, protease), and regulatory proteins. The virus integrates into the host genome as a provirus and can establish latent reservoirs. Without treatment, HIV infection leads to progressive CD4+ T cell depletion, increasing susceptibility to opportunistic infections and malignancies, and ultimately resulting in AIDS and death.

\medterm{Human microbiome}-- The collection of microorganisms that live on and in the human body. The microbiome includes bacteria, viruses, fungi, and archaea. The gut microbiome is the most studied and plays a role in digestion, immune function, and disease. Dysbiosis (alteration of the microbiome) is associated with inflammatory bowel disease, obesity, and infections. Understanding the microbiome is essential for developing new therapeutic approaches including fecal microbiota transplantation and probiotics.

\medterm{Immunodeficiency and infections}-- Immunocompromised individuals are at increased risk for opportunistic infections. Primary immunodeficiencies are genetic disorders affecting immune function. Secondary immunodeficiencies include HIV/AIDS, organ transplantation, chemotherapy, and corticosteroid use. Common opportunistic infections include Pneumocystis jirovecii, CMV, Aspergillus, and Cryptococcus. Prevention through antimicrobial prophylaxis and vaccination is essential. Early diagnosis and treatment of opportunistic infections is critical for improving outcomes in immunocompromised patients, and prophylactic antimicrobial therapy should be tailored to the specific immune defect.

\medterm{Influenza virus}-- A segmented, negative-sense RNA virus belonging to the Orthomyxoviridae family. Influenza causes seasonal epidemics and occasional pandemics of respiratory illness. The virus has two surface glycoproteins: hemagglutinin (HA), which mediates attachment, and neuraminidase (NA), which facilitates release. Influenza A infects humans, birds, and other animals, while influenza B infects only humans. Antigenic drift (point mutations in HA and NA) causes seasonal epidemics, while antigenic shift (reassortment of genome segments) can generate novel pandemic strains to which the population has little pre-existing immunity.

\medterm{Innate immune evasion}-- Mechanisms by which pathogens evade the innate immune system. Strategies include capsules that prevent phagocytosis, biofilms that protect against immune cells, and intracellular survival within macrophages. Some bacteria produce enzymes that degrade complement components. Others modify their surface molecules to avoid recognition. Understanding immune evasion is essential for developing new therapeutic approaches and vaccines. Pathogens that evade innate immunity can establish infections before adaptive immune responses are mounted, making early recognition and rapid diagnosis critical for effective treatment.

\medterm{Integrons}-- Genetic elements that capture and express gene cassettes through site-specific recombination. Integrons contain an integrase gene and a recombination site for gene cassette insertion. They are important vehicles for the spread of antibiotic resistance genes, particularly in gram-negative bacteria. Class 1 integrons are the most clinically significant and are associated with multidrug resistance. Gene cassettes can contain multiple resistance genes, allowing accumulation of resistance. Integrons play a key role in the dissemination of antibiotic resistance in clinical settings, particularly among gram-negative pathogens where they contribute to the emergence of multidrug-resistant strains.

\medterm{Intracellular pathogens}-- Microorganisms that survive and replicate within host cells. Intracellular bacteria include Mycobacterium, Listeria, and Chlamydia. Viruses are obligate intracellular pathogens. Intracellular survival protects pathogens from antibiotics and immune responses. Strategies for intracellular survival include avoiding lysosomal fusion and surviving within phagosomes. Understanding intracellular pathogenesis is essential for developing new therapeutic approaches. Intracellular pathogens require treatment with antimicrobial agents that can penetrate host cell membranes, such as macrolides, tetracyclines, and fluoroquinolones, and often require prolonged therapy to achieve eradication.

\medterm{Legionella pneumophila}-- A gram-negative rod that causes Legionnaires' disease, a severe pneumonia, and Pontiac fever, a milder illness. The organism is found in aquatic environments and is transmitted through inhalation of contaminated water droplets. Risk factors include age, smoking, and immunosuppression. Diagnosis involves urine antigen testing and culture on BCYE agar. Treatment includes fluoroquinolones or azithromycin. Prevention through water system management is essential. Legionella can colonise building water systems, including cooling towers, hot water tanks, and decorative fountains, and outbreaks require identification and decontamination of the environmental source.

\medterm{Leishmania donovani}-- A protozoan parasite that causes visceral leishmaniasis (kala-azar), transmitted by the bite of infected female phlebotomine sandflies. The organism exists as flagellated promastigotes in the sandfly gut and as non-flagellated amastigotes within human macrophages. After inoculation, promastigotes are phagocytosed by macrophages and transform into amastigotes, which multiply intracellularly and disseminate throughout the reticuloendothelial system. Visceral leishmaniasis presents with fever, hepatosplenomegaly, pancytopenia, weight loss, and hypergammaglobulinaemia; it is fatal if untreated. Cutaneous leishmaniasis is caused by other Leishmania species (L. major, L. tropica, L. mexicana), presenting as skin ulcers. Diagnosis involves microscopy of tissue smears (amastigotes in macrophages), culture, and PCR. Treatment includes liposomal amphotericin B and miltefosine.

\medterm{Listeria monocytogenes}-- A gram-positive, facultatively anaerobic rod that causes listeriosis, a foodborne infection. The organism is transmitted through contaminated food, particularly deli meats, soft cheeses, and ready-to-eat foods. It can grow at refrigeration temperatures. Listeriosis is most severe in pregnant women, neonates, and immunocompromised patients. The organism crosses the intestinal barrier and can cross the placenta. Diagnosis involves culture and molecular methods. Treatment includes ampicillin or trimethoprim-sulfamethoxazole, and prompt therapy is critical given the high mortality rate of invasive listeriosis, particularly in neonates and immunocompromised patients.

\medterm{Measles virus}-- A negative-sense RNA virus that causes measles, a highly contagious respiratory illness. The virus is transmitted by respiratory droplets and is one of the most infectious pathogens known. Symptoms include fever, cough, coryza, conjunctivitis, and rash. Complications include pneumonia, encephalitis, and subacute sclerosing panencephalitis (SSPE). Diagnosis involves serology and PCR. Treatment is supportive. Prevention through vaccination with MMR vaccine is highly effective. Measles elimination remains a global health goal, but outbreaks continue to occur in communities with suboptimal vaccination coverage, highlighting the importance of maintaining high immunisation rates.

\medterm{Mobile genetic elements}-- Genetic elements that can move within and between genomes. These include plasmids, transposons, and integrons. Mobile genetic elements can carry antibiotic resistance genes and virulence factors. They facilitate horizontal gene transfer and contribute to the rapid evolution of bacteria. Understanding mobile genetic elements is essential for tracking the spread of resistance and developing new antimicrobial strategies.

\medterm{MRSA}-- Methicillin-resistant Staphylococcus aureus, resistant to beta-lactam antibiotics through acquisition of the mecA gene encoding PBP2a. MRSA is a major cause of healthcare-associated and community-associated infections. Treatment options include vancomycin, daptomycin, and linezolid. MRSA infections are associated with increased morbidity and mortality. Infection prevention includes hand hygiene, contact precautions, and decolonization protocols. The prevalence of MRSA varies by healthcare setting and geographic region, with some areas reporting endemic levels requiring ongoing surveillance and control efforts.

\medterm{MRSA decolonization}-- Strategies to eliminate MRSA carriage from the nares and skin of colonized individuals. Decolonization reduces the risk of infection and transmission. The most common regimen is intranasal mupirocin for 5 days. Chlorhexidine bathing is used for body decolonization. Decolonization is recommended for preoperative patients, those with recurrent infections, and during outbreaks. Eradication is not always successful, and recolonization can occur. Understanding decolonisation strategies is essential for reducing MRSA infection rates in high-risk populations and for preventing surgical site infections in colonised patients.

\medterm{Mucormycosis}-- A life-threatening fungal infection caused by Mucorales order fungi including Rhizopus, Mucor, and Rhizomucor species. The infection primarily affects immunocompromised patients, particularly those with uncontrolled diabetes mellitus and hematologic malignancies. The fungi invade blood vessels causing thrombosis and tissue infarction. Rhino-orbital-cerebral mucormycosis is the most common form. Diagnosis involves culture and microscopy showing broad, ribbon-like, pauciseptate hyphae. Treatment involves surgical debridement of necrotic tissue and antifungal therapy with amphotericin B or posaconazole, with early intervention critical for survival given the high mortality rate.

\medterm{Mumps virus}-- A paramyxovirus that causes mumps, a contagious respiratory illness. The virus is transmitted by respiratory droplets. Symptoms include fever, headache, and parotitis. Complications include orchitis, oophoritis, meningitis, and deafness. Diagnosis involves serology and PCR. Treatment is supportive. Prevention through vaccination with MMR vaccine is highly effective. Mumps outbreaks can occur in vaccinated populations, particularly in close-contact settings.

\medterm{Mycobacterium avium complex}-- A group of slowly growing mycobacteria that cause pulmonary disease and disseminated infection in immunocompromised patients. MAC is the most common nontuberculous mycobacterium causing pulmonary disease. Treatment for pulmonary MAC requires a macrolide (azithromycin or clarithromycin) combined with ethambutol and a rifamycin. Disseminated MAC in HIV patients requires macrolide therapy with antiretroviral therapy. The organism is ubiquitous in the environment. Diagnosis requires isolation of MAC from respiratory or tissue specimens, with molecular methods providing faster identification and susceptibility testing guiding treatment selection.

\medterm{Mycobacterium leprae}-- An acid-fast bacillus that causes leprosy (Hansen's disease), a chronic infection affecting the skin and peripheral nerves. The organism is transmitted through respiratory droplets, though the exact mechanism is not fully understood. The incubation period is long (2-5 years). Tuberculoid leprosy presents with few, well-defined skin lesions. Lepromatous leprosy presents with widespread skin involvement and nerve damage. Diagnosis involves skin biopsy showing acid-fast bacilli. Treatment includes a multi-drug regimen of dapsone, rifampicin, and clofazimine, with therapy duration ranging from 6 months for paucibacillary disease to 12 months or longer for multibacillary disease.

\medterm{Mycobacterium tuberculosis}-- An acid-fast bacillus that causes tuberculosis, a chronic respiratory infection. The organism is transmitted by respiratory droplets and primarily affects the lungs but can disseminate to other organs. Latent TB infection is asymptomatic and non-transmissible. Active TB presents with cough, haemoptysis, fever, and weight loss. Diagnosis involves sputum smear for acid-fast bacilli, culture, and nucleic acid amplification testing. Treatment requires multi-drug therapy for 6-9 months. Isoniazid, rifampicin, pyrazinamide, and ethambutol are used in the intensive phase, followed by isoniazid and rifampicin in the continuation phase, though the emergence of multidrug-resistant TB requires alternative regimens.

\medterm{Mycology}-- The branch of microbiology that studies fungi, including yeasts, moulds, and mushrooms. Medical mycology focuses on fungal pathogens that cause disease in humans. Fungal infections (mycoses) are classified as superficial, cutaneous, subcutaneous, or systemic. Immunocompromised individuals are at increased risk for invasive fungal infections. Common fungal pathogens include Candida, Aspergillus, Cryptococcus, and endemic fungi. Diagnosis involves culture, microscopy, antigen detection, and molecular methods, with rapid molecular tests improving time to diagnosis and facilitating early targeted therapy.

\medterm{Neisseria gonorrhoeae}-- A gram-negative diplococcus that causes gonorrhea, a sexually transmitted infection. The organism has developed resistance to many antibiotics including fluoroquinolones and macrolides. Current treatment recommendations include dual therapy with ceftriaxone and azithromycin, though monotherapy with ceftriaxone is now preferred in some guidelines. The organism colonises mucous membranes of the genitourinary tract, pharynx, and rectum. Complications include pelvic inflammatory disease, infertility, ectopic pregnancy, and disseminated infection, and antimicrobial resistance continues to threaten treatment efficacy.

\medterm{Neisseria meningitidis}-- A gram-negative diplococcus that causes meningitis and meningococcemia. The organism has a polysaccharide capsule that is the target of conjugate vaccines. Serogroups A, B, C, W, X, and Y cause most disease. Meningitis presents with fever, headache, neck stiffness, and altered mental status. Meningococcemia presents with petechial rash and septic shock. Diagnosis involves blood culture and PCR. Treatment includes ceftriaxone or penicillin. Close contacts require prophylaxis with rifampicin, ciprofloxacin, or ceftriaxone to prevent secondary cases following exposure to invasive meningococcal disease.

\medterm{Norovirus}-- A single-stranded RNA virus that is the leading cause of gastroenteritis outbreaks worldwide. The virus is highly contagious and can cause large outbreaks in closed settings including hospitals, cruise ships, and nursing homes. Norovirus causes acute onset of vomiting and diarrhea. Diagnosis is made by PCR or antigen detection in stool. Treatment is supportive with rehydration. The virus is highly stable in the environment and resistant to many disinfectants. There is no vaccine or specific antiviral treatment for norovirus, and management focuses on preventing dehydration through oral or intravenous rehydration.

\medterm{Onchocerca volvulus}-- A filarial nematode that causes onchocerciasis (river blindness), transmitted by repeated bites of infected blackflies (Simulium species) breeding near fast-flowing rivers. Adult worms reside in subcutaneous nodules (onchocercomata), where they can survive for years, producing microfilariae that migrate through the skin and eyes. The host inflammatory response to dying microfilariae causes the pathological manifestations: intense pruritus, dermatitis, depigmentation (leopard skin), lichenification, and, most importantly, progressive ocular damage leading to blindness (keratitis, iridocyclitis, chorioretinitis, optic atrophy). Diagnosis involves skin snip biopsy (microfilariae emerge when skin is placed in saline), slit-lamp examination of the cornea, and PCR. Treatment with ivermectin (annual or semi-annual mass drug administration) kills microfilariae but not adult worms. No safe, effective macrofilaricide currently exists.

\medterm{Parasitology}-- The branch of microbiology that studies parasites, including protozoa, helminths, and ectoparasites. Parasitic infections are a major global health problem, particularly in tropical and subtropical regions. Protozoan parasites include Plasmodium (malaria), Trypanosoma (sleeping sickness, Chagas disease), and Leishmania (leishmaniasis). Helminths include roundworms, tapeworms, and flukes. Diagnosis involves microscopy, serology, and molecular methods. Treatment depends on the parasite and includes antiparasitic agents such as metronidazole, praziquantel, and ivermectin, with prevention focusing on vector control, sanitation, and health education.

\medterm{Plasmids}-- Small, circular, extrachromosomal DNA elements that replicate independently of the bacterial chromosome. Plasmids carry non-essential genes that provide selective advantages including antibiotic resistance, virulence factors, and metabolic capabilities. R-plasmids carry antibiotic resistance genes and can spread resistance between bacteria. F-plasmids carry genes for conjugation. Plasmids are transferred between bacteria through conjugation, transformation, and transduction. They are important vehicles for the dissemination of antibiotic resistance in bacterial populations and are key targets for molecular surveillance of resistance mechanisms.

\medterm{Poliovirus}-- An enterovirus that causes poliomyelitis, a paralytic illness. The virus is transmitted by the fecal-oral route. Most infections are asymptomatic, but a small percentage cause acute flaccid paralysis. The virus attacks motor neurons in the spinal cord. Diagnosis involves stool culture and PCR. There is no specific treatment. Prevention through vaccination with inactivated (Salk) or oral (Sabin) polio vaccine is highly effective. Polio has been eliminated from most countries through vaccination efforts, but transmission continues in Afghanistan and Pakistan, and importation into polio-free areas remains a risk.

\medterm{Polymicrobial infections}-- Infections caused by multiple microbial species. Polymicrobial infections are common in wound infections, intra-abdominal infections, and respiratory infections. The interactions between organisms can be synergistic or antagonistic. Understanding polymicrobial infections is essential for selecting appropriate antimicrobial therapy. Culture and identification of all organisms in a polymicrobial infection can be challenging.

\medterm{Post-exposure prophylaxis}-- The administration of antimicrobial agents after exposure to a pathogen to prevent infection. PEP is used for HIV, hepatitis B, hepatitis C, rabies, and other infections. The timing of PEP is critical, with early administration being most effective. Understanding PEP is essential for managing occupational exposures and outbreaks. PEP guidelines vary by pathogen and exposure type.

\medterm{Rabies}-- A fatal viral encephalitis caused by rabies virus, transmitted through the bite of infected animals. After inoculation, the virus travels to the CNS via peripheral nerves. Symptoms include hydrophobia, aerophobia, and progressive neurological deterioration. Diagnosis is made by direct fluorescent antibody testing of brain tissue. Post-exposure prophylaxis with rabies vaccine and immunoglobulin is highly effective if administered promptly. There is no specific treatment once symptoms develop. Prevention through pre-exposure vaccination of at-risk individuals and animal control programmes, including vaccination of domestic dogs, is essential for reducing the global burden of rabies.

\medterm{Respiratory viruses}-- Viruses that cause respiratory tract infections. Major respiratory viruses include influenza, respiratory syncytial virus (RSV), parainfluenza virus, adenovirus, and human metapneumovirus. SARS-CoV-2 causes COVID-19. Symptoms range from mild upper respiratory illness to severe pneumonia. Diagnosis involves PCR, rapid antigen testing, and viral culture. Treatment is supportive with antiviral agents available for influenza and SARS-CoV-2. Prevention through vaccination and hygiene is essential. Respiratory viruses are responsible for significant morbidity and mortality worldwide, particularly in young children, older adults, and immunocompromised individuals.

\medterm{Retrovirus}-- An RNA virus that replicates by reverse transcribing its RNA genome into DNA, which is then integrated into the host cell genome. Family Retroviridae. Key enzyme: reverse transcriptase (RNA-dependent DNA polymerase) converts viral RNA into double-stranded DNA; integrase inserts the viral DNA (provirus) into the host chromosome. The provirus is replicated with the host cell genome and transcribed into new viral RNA and proteins. Genera: Lentivirus (HIV-1, HIV-2, SIV), Deltaretrovirus (HTLV-1, HTLV-2), Gammaretrovirus (murine leukemia virus), and Spumavirus (foamy viruses). Clinical significance: HIV causes AIDS (depletion of CD4+ T cells, opportunistic infections, malignancies). HTLV-1 causes adult T-cell leukemia/lymphoma and tropical spastic paraparesis (HAM/TSP). Antiretroviral therapy targets reverse transcriptase (NRTIs, NNRTIs), protease, integrase (integrase strand transfer inhibitors, INSTIs), and entry/fusion. Retroviruses are also used as gene therapy vectors (lentiviral vectors) due to their ability to stably integrate into the host genome.

\medterm{Rhinovirus}-- A genus of non-enveloped, single-stranded positive-sense RNA viruses in the Picornaviridae family, the most common cause of the common cold. Rhinoviruses have >160 identified serotypes, each with distinct surface proteins (viral capsid proteins VP1--VP4) that determine antigenicity and host immune response. Transmission: respiratory droplets and direct contact (contaminated hands, fomites). Incubation: 2--3 days. Clinical: nasal congestion, rhinorrhoea, sore throat, cough, sneezing, headache, and mild malaise—symptoms peak at 2--3 days and resolve in 7--10 days. Rhinoviruses replicate optimally at 33--35°C (cooler temperature of the nasal passages) and bind to ICAM-1 (intercellular adhesion molecule 1—major receptor for most serotypes) or LDL receptor (minor group). Immunity is serotype-specific and short-lived, explaining recurrent infections throughout life. Complications: acute sinusitis, acute otitis media (especially in children), exacerbation of asthma and COPD, and pneumonia (in immunocompromised patients). Diagnosis is clinical; PCR or viral culture is rarely needed. Treatment: supportive (rest, hydration, decongestants, NSAIDs, antihistamines). No specific antiviral is approved; pleconaril (viral capsid binder) is investigational. Prevention: hand hygiene, avoiding close contact with infected individuals.

\medterm{Rickettsia}-- A genus of obligate intracellular Gram-negative coccobacillary bacteria transmitted by arthropod vectors (ticks, fleas, lice, mites). Rickettsia are the causative agents of the spotted fever group (SFG) and typhus group (TG) rickettsioses. SFG: Rickettsia rickettsii (Rocky Mountain spotted fever—most severe, mortality 20--30\% untreated), R. conorii (Mediterranean spotted fever, Boutonneuse fever), R. akari (rickettsialpox). TG: R. prowazekii (epidemic typhus—louse-borne, associated with war and poverty, mortality 10--60\% untreated), R. typhi (murine/endemic typhus—flea-borne, milder). Pathophysiology: Rickettsia invade vascular endothelial cells, causing systemic vasculitis, increased vascular permeability, and microthrombi. Clinical: fever, severe headache, myalgia, and a characteristic rash (maculopapular, petechial, or eschar at the bite site—tache noire in SFG; centripetal spread in Rocky Mountain spotted fever; macular rash on the trunk spreading centrifugally in typhus). Diagnosis: serology (IFA—gold standard, Weil–Felix test—obsolete), PCR on skin biopsy or blood, immunohistochemistry of skin biopsy. Treatment: doxycycline 100 mg BD for 7--14 days (first-line for all rickettsioses, including children—short course does not cause tooth staining). Chloramphenicol is an alternative (pregnancy, doxycycline allergy). Prevention: tick avoidance, protective clothing, insect repellent (DEET), and prompt tick removal. No vaccine is available for most rickettsioses.

\medterm{Rotavirus}-- A double-stranded RNA virus that is the leading cause of severe gastroenteritis in children worldwide. The virus has a segmented genome and infects enterocytes of the small intestine. Rotavirus causes profuse watery diarrhea, vomiting, and dehydration. Diagnosis is made by antigen detection in stool. Treatment is supportive with oral rehydration. Vaccination is available and recommended for infants. The virus is highly stable in the environment and transmitted through the fecal-oral route, with vaccination being the most effective preventive measure against severe disease.

\medterm{Rubella virus}-- A togavirus that causes rubella (German measles), a mild febrile illness with rash. The virus is transmitted by respiratory droplets. Rubella is most dangerous during pregnancy, causing congenital rubella syndrome (CRS) with deafness, cataracts, and cardiac defects. Diagnosis involves serology. Treatment is supportive. Prevention through vaccination with MMR vaccine is highly effective. Rubella has been eliminated from the Americas through vaccination. Congenital rubella syndrome is a devastating but preventable consequence of maternal rubella infection, and maintaining high vaccination coverage is essential for elimination.

\medterm{Salmonella}-- See Salmonella in Infectious Diseases.

\medterm{Sarcoptes scabiei}-- A mite that causes scabies, a contagious skin infestation transmitted by prolonged skin-to-skin contact. The female mite burrows into the stratum corneum, laying eggs along the burrow, which appears as a greyish, serpiginous line. The life cycle from egg to adult takes approximately 10-14 days. The host immune response to the mites, their eggs, and their feces (scybala) causes intense pruritus, particularly at night. The typical distribution includes interdigital finger webs, wrists, elbows, axillae, waist, and genitalia, sparing the face in adults. Crusted (Norwegian) scabies is a hyperinfestation form occurring in immunocompromised patients, characterized by thick crusted plaques containing thousands of mites. Diagnosis is clinical, confirmed by microscopy of skin scrapings from burrows (identification of mites, eggs, or scybala). Treatment involves permethrin 5\% cream or oral ivermectin.

\medterm{SARS-CoV-2}-- A coronavirus belonging to the Coronaviridae family that causes COVID-19. SARS-CoV-2 has a positive-sense single-stranded RNA genome and a spike glycoprotein that mediates attachment to ACE2 receptors. The virus replicates in the cytoplasm using its own RNA-dependent RNA polymerase. Transmission occurs through respiratory droplets and aerosols. The disease ranges from mild respiratory symptoms to severe pneumonia and multi-organ failure. Vaccines targeting the spike protein have been developed and deployed globally, with mRNA, viral vector, and protein subunit platforms demonstrating high efficacy against severe disease. The rapid emergence of variants of concern necessitates ongoing vaccine updates and booster strategies.

\medterm{Staphylococcus aureus}-- A gram-positive, catalase-positive, coagulase-positive coccus that causes a wide range of infections including skin infections, pneumonia, endocarditis, and osteomyelitis. S. aureus produces numerous virulence factors including toxins, enzymes, and surface proteins. MRSA strains carry the mecA gene encoding PBP2a. The organism is part of the normal nasal flora. Diagnosis involves culture and coagulase testing. Treatment depends on susceptibility and includes beta-lactams for methicillin-susceptible strains and vancomycin, daptomycin, or linezolid for MRSA, with the choice guided by local epidemiology and patient factors.

\medterm{Streptococcus agalactiae}-- A gram-positive, beta-hemolytic streptococcus (Group B Streptococcus) that is a leading cause of neonatal sepsis, meningitis, and pneumonia. The organism colonises the gastrointestinal and genital tracts. Transmission occurs during birth. Risk factors include prematurity, prolonged rupture of membranes, and maternal colonisation. Screening and intrapartum antibiotic prophylaxis have reduced the incidence of early-onset GBS disease. The organism is identified by Lancefield Group B antigen and hippurate hydrolysis testing, and intrapartum antibiotic prophylaxis has been highly effective in preventing vertical transmission during labour.

\medterm{Streptococcus pyogenes}-- A gram-positive, beta-hemolytic streptococcus (Group A Streptococcus) that causes pharyngitis, skin infections, and post-infectious complications including rheumatic fever and glomerulonephritis. The organism produces streptolysins, erythrogenic toxin, and streptokinase. It is identified by bacitracin sensitivity and Lancefield Group A antigen. Pharyngitis presents with fever, sore throat, and exudates. Skin infections include impetigo and cellulitis. Severe invasive infections include necrotising fasciitis and streptococcal toxic shock syndrome, which require prompt recognition, aggressive surgical debridement, and appropriate antimicrobial therapy.

\medterm{Toxoplasma gondii}-- An obligate intracellular protozoan parasite that causes toxoplasmosis. The definitive host is the cat (family Felidae), which sheds environmentally resistant oocysts in feces. Humans are intermediate hosts, acquiring infection through ingestion of undercooked meat containing tissue cysts, ingestion of oocysts from contaminated soil or cat litter, or transplacentally from mother to fetus. In immunocompetent individuals, infection is usually asymptomatic or causes a mild flu-like illness with lymphadenopathy. Immunocompromised patients (HIV/AIDS, transplant recipients) are at risk of severe disease, most commonly toxoplasma encephalitis presenting with brain abscesses. Congenital toxoplasmosis can cause chorioretinitis, intracranial calcifications, hydrocephalus, and developmental delay. Diagnosis relies on serology (IgM, IgG) and PCR. Treatment involves pyrimethamine with sulfadiazine and folinic acid.

\medterm{Transposons}-- Mobile genetic elements that can move between different locations within the bacterial genome or between the chromosome and plasmids. Transposons carry genes including antibiotic resistance genes and can facilitate the spread of resistance. They are sometimes called "jumping genes." Transposons insert into DNA through a cut-and-paste or copy-and-paste mechanism. They can inactivate genes by inserting into coding regions. Composite transposons consist of insertion sequences flanking antibiotic resistance genes and facilitate their movement between replicons, contributing to the assembly of multidrug-resistant plasmids.

\medterm{Travel medicine}-- The branch of medicine dealing with health problems associated with international travel. Pre-travel consultation includes vaccinations, malaria prophylaxis, and travel advice. Common travel-related illnesses include traveller's diarrhea, malaria, and hepatitis A. High-risk destinations require specific vaccinations including yellow fever, typhoid, and Japanese encephalitis. Post-travel evaluation should consider tropical infections in returning travellers. Understanding travel medicine is essential for providing appropriate pre-travel advice, recognising imported diseases, and preventing the international spread of infectious diseases.

\medterm{Treponema pallidum}-- A spirochete that causes syphilis, a sexually transmitted infection. The organism is transmitted through sexual contact or transplacentally (congenital syphilis). Primary syphilis presents with a painless genital ulcer (chancre). Secondary syphilis presents with rash, mucous patches, and condylomata lata. Tertiary syphilis can involve the cardiovascular system, CNS, and skin. Diagnosis involves serology (RPR, FTA-ABS) and darkfield microscopy. Treatment includes penicillin G. Congenital syphilis remains a preventable cause of stillbirth and neonatal morbidity, and routine antenatal screening and treatment of maternal infection are essential for prevention.

\medterm{Trypanosoma brucei}-- A protozoan parasite that causes African trypanosomiasis (sleeping sickness), transmitted by the bite of the tsetse fly (Glossina species). Two subspecies cause human disease: T. b. gambiense (West African form, chronic, accounting for over 95\% of reported cases) and T. b. rhodesiense (East African form, acute). The parasite is inoculated during a tsetse fly blood meal, multiplies in the blood and lymphatics (haemolymphatic stage), and eventually crosses the blood-brain barrier to invade the central nervous system (meningoencephalitic stage). The haemolymphatic stage presents with intermittent fever, headache, lymphadenopathy (classic Winterbottom sign -- enlarged posterior cervical nodes), and chancre at the bite site. The meningoencephalitic stage causes progressive neurological deterioration, daytime somnolence with night-time insomnia, ataxia, and coma. Diagnosis involves microscopy of blood and CSF, serology, and PCR. Treatment depends on stage: pentamidine or suramin for early disease; eflornithine or melarsoprol for late-stage CNS involvement.

\medterm{Varicella-zoster virus}-- A herpesvirus that causes chickenpox (varicella) and shingles (herpes zoster). VZV establishes latent infections in dorsal root ganglia and can reactivate decades later. Chickenpox is highly contagious and typically occurs in childhood. Reactivation as shingles causes painful dermatomal vesicular rash. Complications include postherpetic neuralgia and disseminated infection in immunocompromised patients. Antiviral treatment with acyclovir or valacyclovir is effective. Vaccination is available for varicella (chickenpox) and herpes zoster (shingles), with live attenuated vaccines recommended for immunocompetent individuals and recombinant zoster vaccine recommended for adults aged 50 years and older.

\medterm{Vector-borne diseases}-- Diseases transmitted by arthropod vectors including mosquitoes, ticks, fleas, and sandflies. Major vector-borne diseases include malaria, dengue, Zika, chikungunya, Lyme disease, and leishmaniasis. Prevention through vector control, personal protection, and chemoprophylaxis is essential. Climate change is affecting the distribution of vectors and vector-borne diseases. Understanding vector biology is essential for developing control strategies. Vector-borne diseases are a major global health problem, accounting for more than 17\% of all infectious diseases and causing over 700,000 deaths annually worldwide, with the burden falling disproportionately on tropical and subtropical regions.

\medterm{Vibrio cholerae}-- A gram-negative, curved rod that causes cholera, a severe diarrhoeal illness. The organism produces cholera toxin, which causes profuse watery diarrhea (rice-water stools). Transmission occurs through contaminated water and food. The infectious dose varies depending on stomach acidity. Diagnosis involves stool culture on selective media (TCBS agar). Treatment is aggressive fluid replacement with oral rehydration solution or intravenous fluids. Antibiotics (doxycycline, azithromycin) reduce the duration and volume of diarrhea and should be given in moderate to severe cases alongside aggressive fluid resuscitation.

\medterm{Viral assembly}-- The process of packaging newly synthesised viral genomes and proteins into complete virions. Assembly typically occurs in the nucleus or cytoplasm, depending on the virus type. The process involves self-assembly of capsid proteins around the viral genome. Enveloped viruses acquire their envelope from host cell membranes during assembly or release. The efficiency of assembly varies by virus type and can be affected by cellular factors. Assembly is a potential target for antiviral drugs. Understanding viral assembly mechanisms provides opportunities for developing antiviral drugs that disrupt capsid formation or packaging of the viral genome.

\medterm{Viral attachment}-- The initial step in viral infection where viral surface proteins bind to specific receptors on host cells. Attachment determines viral tropism and host range. The viral attachment protein is called the ligand, and the host cell receptor is typically a protein or glycan. Examples include HIV gp120 binding to CD4 receptors, influenza hemagglutinin binding to sialic acid, and SARS-CoV-2 spike protein binding to ACE2 receptors. Attachment can be blocked by neutralising antibodies or receptor antagonists, making viral attachment a key target for vaccines and antiviral therapies.

\medterm{Viral culture}-- The growth of viruses in cell culture for identification and characterisation. Viral culture remains the gold standard for many viral infections. The technique involves inoculating cell lines with clinical specimens and observing for cytopathic effects. Shell vial culture uses centrifugation to enhance viral adsorption and provides faster results. Viral culture is more sensitive than antigen detection but less sensitive than PCR. Culture is essential for antiviral susceptibility testing. Understanding viral culture techniques remains important for virology laboratories, particularly for antiviral susceptibility testing and the isolation of novel viruses.

\medterm{Viral cytopathic effects}-- Observable changes in host cells caused by viral infection. CPE include cell rounding, detachment, syncytia formation (fusion of infected cells), inclusion bodies, and cell lysis. CPE are used in viral diagnosis and can be observed by light microscopy. Some viruses cause rapid cell lysis while others cause chronic infection without obvious CPE. The mechanism of CPE varies by virus type and can involve viral protein expression, host cell shutoff, or immune-mediated damage. CPE are important for viral diagnosis and quantification, and the pattern of CPE can provide clues to the identity of the infecting virus.

\medterm{Viral entry}-- The process by which viruses cross the host cell membrane to deliver their genome into the cytoplasm. Entry occurs through membrane fusion (enveloped viruses) or endocytosis (enveloped and non-enveloped viruses). For membrane fusion, viral fusion proteins mediate merger of the viral envelope with the host cell membrane. For endocytosis, the virus is internalised in vesicles and then escapes into the cytoplasm. The mechanism of entry varies by virus type. Entry can be blocked by fusion inhibitors such as enfuvirtide for HIV, making viral entry an important target for antiviral drug development.

\medterm{Viral gastroenteritis}-- Gastrointestinal inflammation caused by viral infections. The major causes include rotavirus, norovirus, adenovirus, and astrovirus. Symptoms include watery diarrhea, vomiting, and abdominal cramps. Diagnosis is made by antigen detection or PCR in stool. Treatment is supportive with rehydration. Prevention through vaccination (rotavirus) and hygiene is essential. Viral gastroenteritis is the most common cause of diarrhoeal illness worldwide. Most cases are self-limited but can cause severe dehydration in young children, older adults, and immunocompromised patients, requiring prompt oral or intravenous rehydration therapy.

\medterm{Viral hepatitis}-- Inflammation of the liver caused by viral infections. The major hepatitis viruses include A, B, C, D, and E. Hepatitis A and E are transmitted through the fecal-oral route and cause acute, self-limited infections. Hepatitis B and C are transmitted through blood and sexual contact and can cause chronic infections leading to cirrhosis and hepatocellular carcinoma. Hepatitis D requires hepatitis B for replication. Diagnosis involves serology and viral load testing. Prevention through vaccination is available for hepatitis A and B, and antiviral therapy can effectively control hepatitis B and C, though access to treatment remains limited in many resource-constrained settings.

\medterm{Viral immune evasion}-- Mechanisms by which viruses evade host immune responses. Strategies include antigenic variation, inhibition of interferon signalling, downregulation of MHC molecules, and encoding viral homologues of cytokines. Understanding viral immune evasion is essential for developing effective vaccines and antiviral therapies. Viruses that evade immune responses can establish persistent infections.

\medterm{Viral latency}-- A state in which viral genomes persist in host cells without producing infectious particles. Latent viruses can reactivate under certain conditions including immunosuppression. Herpesviruses establish latent infections in specific cell types (HSV in neurons, EBV in B cells, CMV in myeloid cells). HIV integrates into the host genome as a provirus and can establish latent reservoirs. Latency is a major barrier to viral eradication. Reactivation can cause disease years after initial infection. Understanding viral latency is critical for developing curative therapies for chronic viral infections, as eradication of latent reservoirs remains one of the greatest challenges in virology.

\medterm{Viral latency mechanisms}-- The strategies by which viruses establish and maintain latent infections. Herpesviruses maintain their genomes as episomes in the nucleus. HIV integrates into the host genome as a provirus. Latent viruses express minimal proteins to avoid immune detection. Reactivation can be triggered by immunosuppression, stress, or other stimuli. Understanding latency mechanisms is essential for developing strategies to cure chronic viral infections. Reactivation can cause disease years after initial infection, and novel therapeutic approaches including latency-reversing agents are being investigated to eliminate persistent viral reservoirs.

\medterm{Viral pathogenesis}-- The mechanisms by which viruses cause disease. Pathogenesis involves attachment, entry, replication, and spread within the host. Viral damage can be direct (cytopathic effects) or indirect (immune-mediated). Direct effects include cell death, altered cell function, and cell transformation. Indirect effects include inflammation, autoimmunity, and cytokine storm. The incubation period, viral load, and host immune response determine disease severity. Some viruses cause acute infections (influenza, norovirus) that are cleared by the immune system, while others establish persistent or latent infections (HIV, HBV, herpesviruses) that require long-term management.

\medterm{Viral release}-- The process by which newly assembled virions exit the host cell to infect other cells. Release occurs through budding (enveloped viruses) or cell lysis (non-enveloped viruses). Budding allows the virus to acquire an envelope from host cell membranes without killing the cell. Cell lysis releases large numbers of virions but destroys the host cell. Neuraminidase inhibitors block release of influenza viruses. Understanding release mechanisms is important for developing antiviral drugs that target the final stages of the viral life cycle, such as neuraminidase inhibitors for influenza and HIV protease inhibitors.

\medterm{Viral replication}-- The process of synthesising new viral genomes and proteins within host cells. The mechanism varies by virus type. DNA viruses typically replicate in the nucleus using host cell DNA polymerases (except poxviruses). RNA viruses typically replicate in the cytoplasm using viral RNA-dependent RNA polymerase. Retroviruses use reverse transcriptase to convert RNA to DNA, which integrates into the host genome. The replication strategy determines susceptibility to antiviral drugs. Viral replication requires exploitation of host cell machinery, and the specific replication strategy of each virus determines its susceptibility to antiviral agents and its capacity for genetic variation.

\medterm{Viral replication cycle}-- The series of steps by which viruses reproduce within host cells. The cycle includes attachment to host cell receptors, entry through membrane fusion or endocytosis, uncoating of the viral genome, replication of viral nucleic acids, transcription and translation of viral proteins, assembly of new virions, and release by budding or cell lysis. The replication cycle varies by virus type and can take hours to days. Some viruses integrate into the host genome (proviruses) and can establish latent infections that persist for the lifetime of the host, requiring lifelong suppressive therapy.

\medterm{Viral structure}-- The basic components of viral particles (virions). Viruses consist of a nucleic acid genome (DNA or RNA) surrounded by a protein coat called a capsid. Some viruses have an outer lipid envelope derived from host cell membranes. The capsid is composed of structural units called capsomeres. Viral genomes can be single-stranded or double-stranded, linear or circular. Enveloped viruses include influenza, HIV, and hepatitis B. Non-enveloped viruses include adenoviruses, noroviruses, and papillomaviruses, and are generally more resistant to environmental inactivation and disinfectants than enveloped viruses.

\medterm{Viral uncoating}-- The process by which the viral capsid is disassembled to release the viral genome into the host cell cytoplasm. Uncoating occurs after entry and before replication. The mechanism varies by virus type and can involve pH-dependent conformational changes, proteolytic cleavage, or disassembly of capsid proteins. For some viruses, uncoating is triggered by the acidic environment of endosomes. For others, uncoating occurs at the cell surface or in the cytoplasm. Uncoating is essential for viral genome release and is a potential target for antiviral drugs that prevent capsid disassembly.

\medterm{VRE}-- Vancomycin-resistant Enterococcus, resistant to vancomycin through acquisition of vanA or vanB gene clusters. VRE is a significant cause of healthcare-associated infections, particularly in immunocompromised patients. Treatment options include linezolid, daptomycin, and quinupristin-dalfopristin (for E. faecium only). VRE infections are associated with increased mortality and healthcare costs. Infection prevention includes hand hygiene, contact precautions, and environmental cleaning.

\medterm{Yeast}-- Unicellular fungi that reproduce by budding (blastoconidia formation) or fission. The clinically most important yeasts are Candida species (Candida albicans, C. glabrata, C. parapsilosis, C. tropicalis, C. krusei, C. auris) and Cryptococcus species (C. neoformans, C. gattii). Candida albicans is a commensal of the human gastrointestinal tract, oral cavity, and vagina, but causes opportunistic infections in immunocompromised patients: oral thrush (white, curd-like plaques on the oral mucosa that scrape off, leaving an erythematous base), vulvovaginal candidiasis (vaginal itching, burning, thick white discharge), esophageal candidiasis (dysphagia, odynophagia, seen in HIV/AIDS with CD4 <100), cutaneous candidiasis (intertrigo in skin folds), and invasive candidiasis/candidaemia (fever, sepsis in ICU patients, neutropenic patients, post-surgery, associated with central venous catheters, treated with echinocandins — caspofungin, micafungin, anidulafungin — or fluconazole for susceptible strains). Cryptococcus neoformans is an encapsulated yeast that causes cryptococcal meningitis in immunocompromised patients (HIV/AIDS with CD4 <100, organ transplant recipients). Saccharomyces cerevisiae (baker's/brewer's yeast) is rarely pathogenic but can cause fungemia in immunocompromised patients. Malassezia species are lipophilic yeasts associated with seborrhoeic dermatitis, pityriasis versicolor, and folliculitis. Diagnosis: KOH preparation (wet mount of clinical specimen shows budding yeasts with or without pseudohyphae), culture (Saboraud dextrose agar, CHROMagar Candida for species identification), and biochemical tests (germ tube test for C. albicans, API 20C, MALDI-TOF). Treatment: topical (clotrimazole, miconazole, nystatin) and systemic (fluconazole, voriconazole, echinocandins, amphotericin B) antifungals, depending on species, site of infection, and susceptibility.

\medterm{Yersinia pestis}-- A gram-negative rod that causes plague, a severe and potentially fatal infection. The organism is transmitted by flea bites or direct contact with infected animals. Bubonic plague presents with painful lymphadenopathy (buboes). Pneumonic plague involves the lungs and is transmitted by respiratory droplets. Septicemic plague causes widespread dissemination. Diagnosis involves culture and PCR. Treatment includes streptomycin, gentamicin, or doxycycline. Plague is a rare but serious infection, with natural reservoirs in rodent populations in various parts of the world, and it remains a public health concern due to its epidemic potential and use as a bioterrorism agent.

\medterm{Zika virus}-- A positive-sense RNA virus belonging to the Flaviviridae family that causes mild febrile illness. Zika virus is transmitted by Aedes mosquitoes and can also be transmitted sexually. Infection during pregnancy can cause congenital Zika syndrome including microcephaly. Most infections are asymptomatic or mild. Diagnosis is made by PCR or serology. There is no specific treatment or vaccine. Prevention through mosquito control and avoiding travel to endemic areas during pregnancy is important. The virus caused a major epidemic in the Americas in 2015-2016, and surveillance continues to monitor its spread and the risk of congenital infections.

\medterm{Zoonotic infections}-- Infections that are naturally transmitted between animals and humans. Major zoonotic infections include rabies, brucellosis, leptospirosis, and psittacosis. The animal reservoir can be domestic or wild animals. Transmission occurs through direct contact, ingestion, or vectors. Prevention through animal vaccination, food safety, and avoiding animal contact is essential. Zoonotic infections are a significant public health problem worldwide. Understanding zoonotic transmission is essential for preventing emerging infectious diseases, as the majority of newly identified pathogens originate from animal reservoirs.
